% === Begin Label Data ===
% ID: 191
% 难度星级: 
% 标签: 
% 备注: 
% 组卷引用次数: 0
% === End  Label Data ===

\begin{problem}{2018}{G}{江苏卷}{26}{数列}
设 $n \in \mathbf{N}^*$, 对 $1, 2, \cdots, n$ 的一个排列 $i_1, i_2, \cdots, i_n$, 如果当 $s<t$ 时, 有 $i_s > i_t$, 则称 $(i_s, i_t)$ 是排列 $i_1, i_2, \cdots, i_n$ 的一个逆序, 排列 $i_1, i_2, \cdots, i_n$ 的所有逆序的总个数称为其逆序数, 例如: 对 $1, 2, 3$ 的一个排列 $231$, 只有两个逆序 $(2, 1), (3, 1)$, 则排列 $231$ 的逆序数为 2. 记 $f_n(k)$ 为 $1, 2, \cdots, n$ 的所有排列中逆序数为 $k$ 的全部排列的个数.

（1）求 $f_3(2), f_4(2)$ 的值;

（2）求 $f_n(2)$ ($n \ge 5$) 的表达式 (用 $n$ 表示).
\end{problem}

\begin{solutions}
(1) 由题意直接求得 $f_3(2)$ 的值，对 $1, 2, 3$ 的排列，

记 $\mu(abc)$ 为排列 $abc$ 的逆序数，对 $1, 2, 3$ 的所有排列，有
$\mu(123)=0, \mu(132)=1, \mu(213)=1, \mu(231)=2, \mu(312)=2, \mu(321)=3$，

因此 $f_3(0)=1, f_3(1)=f_3(2)=2$.

对 $1, 2, 3, 4$ 的排列，利用已有的 $1, 2, 3$ 的排列，将数字 $4$ 添加进去，$4$ 在新排列中的位置只能是最后三个位置，
因此，$f_4(2)=f_3(2)+f_3(1)+f_3(0)=5$.

(2) 对一般的 $n\ (n \ge 4)$ 的情形，可知逆序数为 $0$ 的排列只有一个，逆序数为 $1$ 的排列只能是将排列 $12\cdots n$ 中的任意相邻两个数字调换位置得到的排列，$f_n(1)=n-1$.

为计算 $f_{n+1}(2)$，当 $1, 2, \cdots, n$ 的排列及其逆序数确定后，将 $n+1$ 添加进原排列，$n+1$ 在新排列中的位置只能是最后三个位置，可得 $f_{n+1}(2)=f_n(2)+f_n(1)+f_n(0)=f_n(2)+n$.

则当 $n \ge 5$ 时，$f_n(2)=[f_n(2)-f_{n-1}(2)]+[f_{n-1}(2)-f_{n-2}(2)]+\cdots+[f_5(2)-f_4(2)]+f_4(2)$
$=(n-1)+(n-2)+\cdots+4+f_4(2)=\dfrac{n^2-n-2}{2}$.

因此，当 $n \ge 5$ 时，$f_n(2)=\dfrac{n^2-n-2}{2}$.
\end{solutions}
